\section{Equal-length words and bounded realization}\label{sec:sync45}
All groups in this section act faithfully on the active space $R$. Choose a reference command $A=U_{a_0}$. For $a\in\calA$ and $t\in\Z$, write
\[
 d_a=A^{-1}U_a,\qquad b_{t,a}=A^{-t}d_a A^t.
\]

\begin{lemma}[Tail generators and reference independence]\label{lem:tail45}
For every integer $T$, the closed subgroup generated by
$\{b_{t,a}:t\ge T,\ a\in\calA\}$ is $\sync$. Equivalently, $\sync$ is generated topologically by all equal-length differences $U_v^{-1}U_w$, $|v|=|w|$. It does not depend on the reference command.
\end{lemma}
\begin{proof}
There are positive integers $n_i\to\infty$ with $A^{n_i}\to I$. For a finite-order matrix take multiples of the order; otherwise this follows by taking arbitrarily close pairs of powers in the compact closure of the cyclic subgroup. Consequently $b_{t+n_i,a}\to b_{t,a}$ for every integer $t$. Thus a closed tail-generated subgroup contains every $b_{t,a}$, $t\in\Z$. The closed subgroup $L$ generated by these elements is invariant under conjugation by $A$ and $A^{-1}$. It contains each $d_a$, and therefore is invariant under conjugation by every $U_a=A d_a$ and its inverse. By continuity it is normal in $H$. Since $U_a^{-1}U_b=d_a^{-1}d_b$, it contains every generator in \eqref{eq:sync-intro}; conversely all the $b_{t,a}$ belong to $\sync$. Hence $L=\sync$.

Every $b_{t,a}$ with $t\ge0$ is an equal-length difference: take the reference word of length $t+1$ and the word whose product is $U_a A^t$. On the other hand, all length-$n$ word products have the same coset $A^n\sync$ because the quotient identifies every command. All their differences lie in $\sync$. These two inclusions prove the equivalent description and the reference independence.
\end{proof}

The recurrence in the first paragraph of the proof can also be verified by simultaneous approximation of the finitely many eigenangles of $A$. No numerical recurrence rate is used.

\begin{lemma}[Stabilization of normalized word orbits]\label{lem:stabilize45}
For a unit vector $u$ in an $H$-invariant subspace, define
\[
 F_t(u)=\{A^{-t}U_wu:w\in\calA^t\},\qquad F_0(u)=\{u\}.
\]
Then $F_t(u)\subseteq F_{t+1}(u)$. Their cardinalities are bounded in $t$ if and only if the orbit $\sync u$ is finite. More precisely, if $|F_t(u)|\le k$ for all $t$, then $|\sync u|\le k$ and $\sync^\circ u=u$.
\end{lemma}
\begin{proof}
Appending the reference command to a word gives the nesting. More generally
\begin{equation}\label{eq:normalized-recursion45}
 F_{t+1}(u)=\bigcup_{a\in\calA}b_{t,a}F_t(u).
\end{equation}
If the cardinalities are bounded, the increasing finite sets stabilize: $F_t(u)=F$ for all $t\ge T$, with $|F|\le k$. Each $b_{t,a}$ maps $F$ into $F$; its invertibility makes that map a permutation. The group generated by the tail therefore preserves $F$. Its closure also preserves $F$, since $F$ is finite and closed. Lemma~\ref{lem:tail45} implies $\sync u\subseteq F$. A connected group acts trivially on a finite orbit, proving the last assertion. Conversely $A^{-t}U_w\in\sync$, so a finite $\sync$-orbit contains every $F_t(u)$.
\end{proof}

\begin{lemma}[A positive packet on the moving subspace]\label{lem:positive-packet45}
Let $E$ be an $H$-invariant real subspace with $E^{\sync^\circ}=\{0\}$. For every $k\ge1$ there is a length $B_k\ge1$ such that every unit $u\in E$ has
\[
 |\{U_wu:w\in\calA^{B_k}\}|>k.
\]
For the uniform probability law $p_k$ on all length-$B_k$ words, the distortion of \eqref{eq:distortion43}, computed in $E$, satisfies
\begin{equation}\label{eq:positive-gamma45}
 \gamma_k:=D_k^E(p_k)>0.
\end{equation}
\end{lemma}
\begin{proof}
The sets
\[
 C_t=\{u\in\Sphere(E):|F_t(u)|\le k\}
\]
are closed: for a finite collection of matrix products, the cardinality condition is a finite union of sets specified by vector equalities. They are nested by Lemma~\ref{lem:stabilize45}. If all were nonempty, compactness of $\Sphere(E)$ would give a unit $u$ in their intersection. The same lemma would give $\sync^\circ u=u$, contrary to the hypothesis. Thus one $C_{B_k}$ is empty. Multiplication by $A^{B_k}$ does not change cardinality.

For a full-support word law, a scalar-product average in \eqref{eq:distortion43} can equal one only if every unit vector $U_wu$ equals one of the $k$ centers. This is excluded for every $u$ by the first assertion. Compactness of the direction and center spaces makes the maximum strictly less than one. This proves \eqref{eq:positive-gamma45}. This argument is qualitative: it supplies no dimension-free lower bound for $\gamma_k$.
\end{proof}

\begin{proof}[Proof of Theorem~\ref{thm:main45}]
Suppose first that $\sync$ is finite. Use labels $(x,\ell)\in X\times\sync$, initialize at $(x,I)$, and at command epoch $t+1$ update deterministically by
\begin{equation}\label{eq:finite-tracker45}
 (x,\ell)\longmapsto(x,b_{t,a}\ell).
\end{equation}
An induction gives $\ell=A^{-t}U_w$ after the first $t$ letters. The decoder at horizon $N$ is
\[
 d_j(x,\ell)=\rho\ip{e_j}{A^N\ell x}\in[-1,1].
\]
It gives \eqref{eq:target} exactly. Only the label persists, and the time-dependent permutation in \eqref{eq:finite-tracker45} uses the already declared external epoch. The held answer adds the maximum with two, not a free output buffer.

Now suppose $\sync$ is infinite. Being a closed subgroup of a compact matrix group, it is a compact Lie group. A compact zero-dimensional Lie group is finite; hence its identity component is nontrivial. Let $E=R\ominus R^{\sync^\circ}$. Normality of $\sync^\circ$ in $H$ makes both summands $H$-invariant and the orthogonal projection $P_E$ commute with all commands. If $E=0$, $\sync^\circ$ would act trivially on the faithful active space, a contradiction. If every seed had zero projection to $E$, its entire $H$-orbit would also, contradicting the definition of $R$. Therefore $\eta_X>0$.

Choose a seed $x$ attaining $\eta_X$ and follow the unit orbit
\[
 Y_t=U_{a_t}\cdots U_{a_1}P_E x/\eta_X\in E.
\]
The actual machine state is not changed or enlarged. For each complete word its decoder mean differs coordinatewise from $\rho U_wx$ by at most $2\varepsilon$. Projection and pairing with $Y_N$ give
\[
 \E\ip{Y_N}{P_Ed(S_N)}\ge\rho\eta_X-2\varepsilon\sqrt D.
\]
Since $\norm{P_Ed}\le\sqrt D$, its conditional-centroid amplitude satisfies
\begin{equation}\label{eq:activated-calibration45}
 q_N\ge\kappa_X:=\rho\eta_X/\sqrt D-2\varepsilon>0.
\end{equation}
This is one analytic pairing of the separate coordinate decoders, not a requirement to answer several queries jointly.

Fix a proposed peak $k$. Apply Lemma~\ref{lem:positive-packet45} on $E$, and use independent complete uniform packets of length $B_k$ followed by a fixed deterministic remainder. Lemma~\ref{lem:packet} applies at their endpoints to the actual hidden register, regardless of its intermediate representations. On the remainder conditional Jensen cannot increase amplitude. Since $q_0=1$, every such machine must satisfy
\begin{equation}\label{eq:threshold45}
 \kappa_X\le(1-\gamma_k)^{\lfloor N/B_k\rfloor}.
\end{equation}
If $\gamma_k=1$, already one complete packet is excluded; otherwise
$\lfloor N/B_k\rfloor[-\log(1-\gamma_k)]\le\log(1/\kappa_X)$.
Thus this fixed $k$ fails for every sufficiently large $N$. The conclusion is a limit to infinity, not merely an unbounded subsequence. Taking $\varepsilon=0$ proves the exact converse.
\end{proof}

\begin{corollary}[A uniform basis-seed threshold]\label{cor:basis45}
For $X=\{\pm e_1,\ldots,\pm e_D\}$ and fixed $\varepsilon<\rho/(2D)$, bounded width is equivalent to finiteness of $\sync$. Otherwise $W_{N,\varepsilon}\to\infty$.
\end{corollary}
\begin{proof}
If $E\ne0$, then $\sum_j\norm{P_Ee_j}^2=\dim E\ge1$, so $\eta_X\ge D^{-1/2}$. Apply Theorem~\ref{thm:main45}.
\end{proof}

\begin{corollary}[Every finite planar rotation alphabet]\label{cor:planar-dichotomy45}
For four coordinate seeds in the plane, commands with angles
$2\pi\alpha_0,\ldots,2\pi\alpha_r$, and fixed
$\varepsilon<\rho/(2\sqrt2)$, width is bounded if and only if every
$\alpha_i-\alpha_0$ is rational. If all differences have common denominator
$q$, an exact machine uses at most $\max(2,4q)$ labels. Otherwise width tends to infinity. The identity letter is not required.
\end{corollary}
\begin{proof}
The commands commute, so $\sync$ is the closed subgroup generated by their difference rotations. It is finite exactly in the rational-difference case. Otherwise it is the whole circle group, $E=\R^2$ and $\eta_X=1$. The common rotation $N\alpha_0$ appears only in the horizon-dependent decoder of \eqref{eq:finite-tracker45}.
\end{proof}

\begin{remark}[Noncommuting examples and unused coordinates]\label{rem:examples45}
Let $A_0=R_\theta\oplus I_3$ on $\R^2\oplus\R^3$, with $\theta/(2\pi)$ irrational, and use $A_0$, $R_\theta\oplus P$, $R_\theta\oplus Q$, where $P,Q$ generate the six permutation matrices on three coordinates. The commands do not commute. Their active group is the infinite group $\SO(2)\times S_3$, but $\sync=\{I_2\}\times S_3$ is finite, giving at most sixty labels for the ten basis seeds. Conversely, if the active group is $\SO(3)$ and two commands differ, its synchronization subgroup is $\SO(3)$, since this group has no nontrivial proper closed normal subgroup. The width therefore diverges without importing a spectral-gap theorem. Finally, rotations on unused coordinates contribute nothing: for seeds on their common fixed axis, restriction to $R$ makes the criterion finite. These examples distinguish synchronized action, ambient generated group, and seed activation.
\end{remark}
